% ---------------------------------------------------------------------------
% Conference paper draft - Finance-AI fraud model.
%
% Compile with the IEEEtran class (IEEE conference format):
%     pdflatex paper.tex
%     pdflatex paper.tex   (twice, for references)
%
% AUTHORS ARE PLACEHOLDERS - replace every bracketed field before submission.
%
% Every number in this paper is reproduced from the repository's own pipeline
% (see Section IX and docs/DATASET.md). No result here is estimated,
% extrapolated, or copied from an unrelated source.
% ---------------------------------------------------------------------------
\documentclass[conference]{IEEEtran}

\usepackage{cite}
\usepackage{amsmath,amssymb,amsfonts}
\usepackage{graphicx}
\usepackage{textcomp}
\usepackage{xcolor}

\begin{document}

\title{Feature-Starved Fraud Detection: An Honest Evaluation of a\\
Transaction-Level Classifier Under a Production Scoring Contract}

\author{
\IEEEauthorblockN{[Author Name]}
\IEEEauthorblockA{[Department]\\
[Institution]\\
[City, Country]\\
[email@example.com]}
\and
\IEEEauthorblockN{[Co-Author Name]}
\IEEEauthorblockA{[Department]\\
[Institution]\\
[City, Country]\\
[email@example.com]}
}

\maketitle

\begin{abstract}
Machine-learned transaction fraud detection is routinely reported with
near-perfect metrics, but those figures usually depend on features---PCA
components, merchant identities, counterparty graphs---that a live scoring
service may not be able to compute. This paper studies the practical
consequence of that gap. We train and compare three classifiers
(logistic regression, random forest, and gradient-boosted trees) on a real,
publicly documented fraud label set (the ULB/Worldline credit-card dataset,
OpenML id 1597), using \emph{only} the features that a production scoring
contract can supply at decision time: amount, cyclical time-of-day, and
coarse categorical flags. The PCA components (\texttt{V1}--\texttt{V28}) that
carry most of the source signal are deliberately excluded because the serving
path cannot reproduce them. We report two evaluations. Under a stratified
hold-out, the best model (random forest) reaches ROC-AUC $0.7577$ and
precision--recall AUC (PR-AUC) $0.0866$, roughly $53\times$ the $0.164\%$
base rate. Under a \emph{chronological} hold-out---the stricter, more
realistic protocol---PR-AUC collapses to $\approx 0.002$, barely above the
base rate, even though ROC-AUC stays near $0.66$--$0.68$. We argue that the
chronological figure is the honest headline, that the model is best deployed
as a weak secondary signal behind a deterministic rule engine, and that
reporting both splits is necessary to avoid overstating fraud-detection
performance. All results are reproducible from the released pipeline, and the
trained artifact is integrity-checked with a SHA-256 digest.
\end{abstract}

\begin{IEEEkeywords}
fraud detection, imbalanced classification, precision--recall, temporal
validation, feature engineering, reproducibility, random forest, XGBoost
\end{IEEEkeywords}

\section{Introduction}
\IEEEPARstart{T}{ransaction} fraud detection is one of the most cited
applications of machine learning in finance, and published studies frequently
report AUCs above $0.95$. Those numbers are usually obtained on
anonymised benchmarks whose strongest predictors are principal components or
aggregate behavioural variables unavailable to the system actually making the
decision. When the same model is moved behind a real scoring API, the
available features shrink, and the reported performance no longer describes
the deployed system.

This paper is a deliberately narrow, honest case study. We ask a single
question: \emph{how well can a transaction-level fraud classifier perform
when restricted to the features a production scoring contract can actually
compute?} We answer it on a real, documented dataset, with two evaluation
protocols, and we report the less flattering protocol as the headline.

Our contributions are:
\begin{itemize}
\item A feature set of 31 interpretable features that are computable at
scoring time from raw transaction fields (amount, timestamp, type, category,
merchant, source), with no future information and no label leakage
(Section~IV).
\item A comparison of three standard classifiers under identical splits and
preprocessing, selected by precision--recall AUC rather than accuracy
(Sections~VI--VII).
\item A stratified-versus-chronological contrast that quantifies how much a
random split flatters the model, and an explicit statement of the deployment
role the model can honestly fill (Sections~VII--VIII).
\item A fully reproducible pipeline whose exported artifact is verified by
SHA-256 digest (Section~IX).
\end{itemize}

We do not claim a production-grade detector. We claim a rigorously measured
\emph{weak} signal, and we show why that measurement is the useful one.

\section{Related Work}
Credit-card fraud detection on highly imbalanced data is well studied. The
dataset we use originates from Dal Pozzolo et al.~\cite{dalpozzolo2015}, who
also discuss undersampling calibration for unbalanced problems. Class
imbalance is central to this literature: resampling methods such as
SMOTE~\cite{chawla2002} are common, and the choice of metric is known to
matter greatly. Davis and Goadrich~\cite{davis2006} established the
relationship between precision--recall and ROC curves, and Saito and
Rehmsmeier~\cite{saito2015} showed that PR curves are more informative than
ROC curves on imbalanced data---the reason we select models by PR-AUC.
Fawcett~\cite{fawcett2006} gives the standard treatment of ROC analysis.
For the learners, we use logistic regression, random
forests~\cite{breiman2001}, and gradient-boosted trees~\cite{chen2016}. At a
higher level, anomaly-detection surveys~\cite{chandola2009} place fraud in the
broader family of rare-event detection, and Sculley et
al.~\cite{sculley2015} describe how train/serve feature skew---the exact
failure mode our parity tests guard against---accumulates real technical debt.
Unlike work that reports only the in-distribution split, we foreground the
temporal split, which better reflects deployment.

\section{System and Scoring Contract}
The detector is one component of an end-to-end personal-finance platform with
a Python/FastAPI backend, a relational store, and web and mobile clients. For
this paper, only the scoring contract matters.

\subsection{The 31-Feature Contract}
Scoring consumes a fixed 31-dimensional vector built from raw transaction
fields by \texttt{app.ml.features.build\_feature\_vector}. The offline
training package (\texttt{finance\_ai\_ml}) reproduces the same 31 features
with the same names. The two implementations intentionally order the columns
differently---the service in semantic order, the offline package
alphabetically---so the exported model artifact records its own feature order
and scoring projects the input vector onto \emph{that} recorded order
($\text{pred}\leftarrow f(x[\text{name}])$ for each recorded name). A rename
on one side only would silently drop a column and degrade every score, so the
repository ships automated parity tests that assert both the \emph{names} and
the \emph{values} agree; a drifted column set is treated as a hard error, not
a warning.

\subsection{Rule Engine and Model Blend}
A deterministic rule engine provides the first line of defence and is always
available. When a trained model is present, the final score is a weighted
blend
\begin{equation}
\text{combined} = \mathrm{round}\!\left(0.7\cdot \text{rule} +
0.3\cdot \text{ml}\right),
\end{equation}
clamped to $[0,100]$, and the rule verdict acts as a floor:
$\text{effective} = \max(\text{combined},\,\text{rule})$. A transaction the
rules flag is therefore never unflagged by the model. This design encodes the
expectation---borne out by Section~VII---that the model is a \emph{secondary}
signal.

\section{Feature Engineering}
All 31 features are computed from fields available at scoring time. They fall
into four families.

\paragraph{Amount (7 features).}
Raw amount; logarithmic and square-root transforms
(\texttt{amount\_log}, \texttt{amount\_sqrt}) to compress the heavy right
tail; and three relative features computed against the transaction owner's own
mean amount: the ratio \texttt{amount\_vs\_mean}, its deviation
\texttt{amount\_deviation} (ratio~$-1$), and a binary
\texttt{is\_amount\_outlier} (ratio $>3$). A binary
\texttt{is\_round\_amount} captures the common ``round-number'' fraud
signature (amount $\ge 100$ and divisible by $100$).

\paragraph{Temporal (14 features).}
The hour is encoded cyclically as \texttt{hour\_sin}/\texttt{hour\_cos} so
that $23{:}00$ and $00{:}00$ are adjacent; \texttt{weekday} is ordinal, with a
seven-way one-hot (\texttt{dow\_0}--\texttt{dow\_6}) and an
\texttt{is\_weekend} flag. A binary \texttt{is\_night} marks
$22{:}00$--$06{:}00$; \texttt{day\_of\_month} and \texttt{month} are included
as-is.

\paragraph{Categorical (8 features).}
One-hot indicators for transaction type
(\texttt{is\_expense}, \texttt{is\_income}) and source
(\texttt{source\_manual}, \texttt{source\_ocr}, \texttt{source\_sms}), an
EMI flag (\texttt{is\_emi}), and presence flags
\texttt{category\_known} and \texttt{merchant\_known}.

\paragraph{Behavioural aggregates (2 features).}
\texttt{merchant\_tx\_count} (how often this user has transacted with this
merchant) and \texttt{user\_tx\_count} (how many transactions this user has).
These are computed over the same static dataset and therefore, on a single
pseudo-user, are near-constant (see Section~VIII).

\paragraph{Leakage controls.}
The train/test split is taken \emph{before} any preprocessing is fitted; the
linear model's \texttt{StandardScaler} lives inside a scikit-learn
\texttt{Pipeline} and is fitted on the training fold only; class imbalance is
handled at the estimator, never by resampling the test set; and
\texttt{amount\_vs\_mean} uses the user's overall mean, matching the serving
path so no future window leaks into a training row.

\section{Dataset}
We use the ULB/Worldline \emph{Credit Card Fraud Detection} dataset distributed
via OpenML (data id 1597)~\cite{vanschoren2014,dalpozzolo2015}. It is a real,
documented, anonymised label set: 284{,}807 transactions, of which 492 are
confirmed fraud ($0.172\%$). Each row carries a timestamp (\texttt{Time}), an
amount (\texttt{Amount}), the label (\texttt{Class}), and 28
privacy-protecting PCA components (\texttt{V1}--\texttt{V28}).

\paragraph{Why the PCA components are dropped.}
\texttt{V1}--\texttt{V28} are not reconstructable into merchant, category, or
counterparty fields, and the production scoring contract cannot produce them.
A model that relied on them would be unscorable against a live transaction.
They are therefore excluded---the deliberate cause of the ``feature
starvation'' analysed below.

\paragraph{Mapping and cleaning.}
\texttt{Time} is mapped to a timestamp ($2013\text{-}09\text{-}01$ UTC $+$
seconds); \texttt{Amount} is unchanged; \texttt{Class} becomes the label.
Every other training field takes its documented default (single synthetic
user, \texttt{expense}, uncategorised). Rows with non-positive or
non-numeric amounts are dropped by the pipeline's standard rule, removing
1{,}825 rows including 27 zero-amount frauds, leaving 282{,}982 rows and 465
positives ($0.1643\%$). Table~\ref{tab:data} summarises the data.

\begin{table}[t]
\centering
\caption{Dataset after acquisition and the positive-amount cleaning rule.}
\label{tab:data}
\begin{tabular}{@{}ll@{}}
\hline
\textbf{Field} & \textbf{Value} \\
\hline
Source & OpenML id 1597 (ULB / Worldline) \\
Raw rows & 284{,}807 ($492$ fraud, $0.172\%$) \\
Modelled rows & 282{,}982 ($465$ positives, $0.1643\%$) \\
Rows dropped (amount $\le 0$) & 1{,}825 (incl.\ 27 frauds) \\
Date span & 2013-09-01 $\to$ 2013-09-02 \\
Amount (min / mean / max) & 0.01 / 88.92 / 25{,}691.16 \\
Exact duplicate rows & 4{,}822 \\
\hline
\end{tabular}
\end{table}

\section{Experimental Setup}
\paragraph{Candidates.}
Three estimators are compared under identical splits and preprocessing:
\begin{itemize}
\item \textbf{Logistic regression (LR)}: \texttt{StandardScaler} followed by
\texttt{LogisticRegression(max\_iter=2000, class\_weight=balanced)} in one
pipeline.
\item \textbf{Random forest (RF)}: 400 trees, \texttt{min\_samples\_leaf=2},
\texttt{class\_weight=balanced\_subsample}.
\item \textbf{Gradient boosting (XGB)}: 200 estimators, depth 4, learning rate
0.05, subsample 0.9, column subsample 0.9, $\lambda=1.0$, with
\texttt{scale\_pos\_weight}$\,=\,(1-p)/p$ to counter the imbalance.
\end{itemize}
The decision threshold is $0.5$ throughout.

\paragraph{Splits.}
A fixed $20\%$ hold-out is used, under two protocols:
\emph{stratified} (random, class-balanced) and \emph{chronological} (oldest
$80\%$ train, newest $20\%$ test, no shuffling). The random seed is fixed
(42), so a rerun reproduces the numbers exactly.

\paragraph{Metrics.}
We report ROC-AUC, PR-AUC (average precision), precision, recall, F1, and the
raw false-positive count. Models are \emph{selected} by PR-AUC---the correct
choice when positives are under $1\%$~\cite{saito2015}.

\section{Results}
\subsection{Stratified hold-out}
The stratified test set contains 56{,}597 rows with 93 positives (prevalence
$0.164\%$). Table~\ref{tab:strat} reports the results. Random forest attains
the best PR-AUC ($0.0866$); XGBoost has the highest ROC-AUC ($0.8278$) but a
materially worse PR-AUC ($0.0544$) at threshold $0.5$, which is why the
selection metric is PR-AUC. At the $0.164\%$ base rate, PR-AUC $0.0866$ is
roughly $53\times$ random ranking. Even so, precision is $2.13\%$: the model
that catches a third of frauds also raises 1{,}424 false positives.

\begin{table}[t]
\centering
\caption{Stratified $20\%$ hold-out (56{,}597 rows, 93 positives).
Selected model in bold.}
\label{tab:strat}
\begin{tabular}{@{}lcccccr@{}}
\hline
\textbf{Model} & \textbf{ROC-AUC} & \textbf{PR-AUC} & \textbf{Prec.} &
\textbf{Rec.} & \textbf{F1} & \textbf{FP} \\
\hline
LR  & 0.6844 & 0.0044 & 0.0031 & 0.6129 & 0.0061 & 18{,}472 \\
\textbf{RF}  & 0.7577 & \textbf{0.0866} & 0.0213 & 0.3333 & 0.0401 & 1{,}424 \\
XGB & 0.8278 & 0.0544 & 0.0065 & 0.6882 & 0.0129 & 9{,}759 \\
\hline
\end{tabular}
\end{table}

\subsection{Chronological hold-out}
The chronological test set contains 56{,}597 rows with 70 positives. As
Table~\ref{tab:chrono} shows, performance collapses. The best PR-AUC is
$0.0023$---barely above the base rate ($70/56{,}597\approx 0.124\%$). Random
forest flags nothing at threshold $0.5$; the others achieve moderate recall
only by accepting thousands of false positives. ROC-AUC nevertheless remains
near $0.66$--$0.68$, which is the tell-tale signature of a classifier that
\emph{ranks} slightly better than chance but cannot separate a rare positive
class with any practical precision.

\begin{table}[t]
\centering
\caption{Chronological $20\%$ hold-out (56{,}597 rows, 70 positives).}
\label{tab:chrono}
\begin{tabular}{@{}lccccr@{}}
\hline
\textbf{Model} & \textbf{ROC-AUC} & \textbf{PR-AUC} & \textbf{Rec.} &
\textbf{F1} & \textbf{FP} \\
\hline
LR  & 0.6538 & 0.0023 & 0.1571 & 0.0032 & 6{,}702 \\
RF  & 0.6626 & 0.0021 & 0.0000 & 0.0000 & 0 \\
XGB & 0.6827 & 0.0022 & 0.2429 & 0.0052 & 6{,}442 \\
\hline
\end{tabular}
\end{table}

\subsection{Interpretation}
The stratified split is flattering partly because a random partition lets
near-duplicate transactions (4{,}822 byte-identical rows exist) fall on both
sides of the split. The chronological split, which forbids that leakage,
is the honest estimate of forward-looking performance. We therefore treat
PR-AUC $\approx 0.002$ as the headline and the $53\times$ lift as an
upper, in-distribution bound. This is consistent with deploying the model at
only $30\%$ weight behind a rule engine (Section~III-B).

\section{Discussion and Limitations}
The modest numbers are explained, not excused, by the following:

\begin{itemize}
\item \textbf{Feature starvation (dominant).} The signal that makes this
dataset separable lives in \texttt{V1}--\texttt{V28}, which the serving
contract cannot supply. Only amount and time-of-day reach the model.
\item \textbf{Distribution shift.} All rows share one synthetic user, so
\texttt{amount\_vs\_mean}, \texttt{merchant\_tx\_count}, and
\texttt{user\_tx\_count} are computed over a single pseudo-user; a real
user's per-transaction baseline differs.
\item \textbf{No merchant/category signal.} The source lacks merchant and
category values, so those one-hot features are constant and carry no weight.
\item \textbf{Zero-amount frauds dropped.} 27 of 492 frauds have amount $0$
and are removed by the positive-amount rule; this bias is unavoidable without
fabricating amounts.
\item \textbf{Two-day window.} \texttt{month} and \texttt{day\_of\_month} take
almost no range, and weekly seasonality is unrepresented.
\end{itemize}

\paragraph{What the pipeline does prove.}
That a real model can be trained, evaluated, selected, exported, and scored
end-to-end on real labels, with integrity checking and train/serve parity
enforced by tests. A production-grade detector would require data that
exposes merchant, category, and counterparty features \emph{at scoring time}.

\section{Reproducibility}
The dataset is public and is not redistributed with the pipeline. It is
acquired and mapped by a documented script, then compared:
\begin{quote}
\small\ttfamily
python scripts/prepare\_ulb\_creditcard.py --work-dir ./ulb\\
python -m finance\_ai\_ml compare ./ulb/ulb\_creditcard.csv\\
\hspace*{1em}--split \{stratified|chronological\}
\end{quote}
The experiments reported here were reproduced with Python 3.14.0, NumPy
2.5.2, pandas 3.0.5, scikit-learn 1.9.0, and XGBoost 3.4.1, and matched the
recorded metrics to four decimal places (including exact false-positive
counts). The exported artifact records its feature order, an artifact-format
version, and a SHA-256 digest of the model binary; the digest was independently
recomputed and matched
(\texttt{734b52d3\ldots a3b24}). Random seeds, thresholds, and the
preprocessing policy are fixed in a single configuration object, so a rerun
reproduces the same model.

\paragraph{Note on the deployed artifact.}
The shipped artifact corresponds to the \emph{stratified} selection
(random forest, PR-AUC $0.0866$). Because the chronological estimate is far
lower, the model is documented and used as a weak secondary signal; the
serving blend and the rule-floor guarantee in Section~III-B are part of that
honest posture.

\section{Conclusion}
Restricting a fraud classifier to the features a production system can
compute at decision time changes the story entirely. On a real, documented
dataset, a random forest reaches PR-AUC $0.0866$ under a stratified hold-out
but only $\approx 0.002$---near the base rate---under a chronological
hold-out, despite ROC-AUC staying near $0.66$--$0.68$. We argue that the
chronological number is the one to report, that PR-AUC (not accuracy or ROC)
should drive selection, and that such a model belongs behind a deterministic
rule engine rather than in front of one. The full pipeline reproduces these
numbers exactly and integrity-checks its output, so the results are auditable
rather than asserted.

\begin{thebibliography}{00}
\bibitem{dalpozzolo2015} A.~Dal Pozzolo, O.~Caelen, R.~A. Johnson, and
G.~Bontempi, ``Calibrating probability with undersampling for unbalanced
classification,'' in \emph{Proc. IEEE Symp. Series on Computational
Intelligence (SSCI)}, 2015.
\bibitem{chawla2002} N.~V. Chawla, K.~W. Bowyer, L.~O. Hall, and
W.~P. Kegelmeyer, ``SMOTE: Synthetic minority over-sampling technique,''
\emph{Journal of Artificial Intelligence Research}, vol.~16, pp.~321--357,
2002.
\bibitem{davis2006} J.~Davis and M.~Goadrich, ``The relationship between
precision--recall and ROC curves,'' in \emph{Proc. Int. Conf. Machine Learning
(ICML)}, 2006.
\bibitem{saito2015} T.~Saito and M.~Rehmsmeier, ``The precision--recall plot
is more informative than the ROC plot when evaluating binary classifiers on
imbalanced datasets,'' \emph{PLoS ONE}, vol.~10, no.~3, 2015.
\bibitem{fawcett2006} T.~Fawcett, ``An introduction to ROC analysis,''
\emph{Pattern Recognition Letters}, vol.~27, no.~8, pp.~861--874, 2006.
\bibitem{breiman2001} L.~Breiman, ``Random forests,'' \emph{Machine
Learning}, vol.~45, no.~1, pp.~5--32, 2001.
\bibitem{chen2016} T.~Chen and C.~Guestrin, ``XGBoost: A scalable tree
boosting system,'' in \emph{Proc. ACM SIGKDD Int. Conf. Knowledge Discovery
and Data Mining (KDD)}, 2016.
\bibitem{chandola2009} V.~Chandola, A.~Banerjee, and V.~Kumar, ``Anomaly
detection: A survey,'' \emph{ACM Computing Surveys}, vol.~41, no.~3, 2009.
\bibitem{sculley2015} D.~Sculley \emph{et al.}, ``Hidden technical debt in
machine learning systems,'' in \emph{Advances in Neural Information Processing
Systems (NeurIPS)}, 2015.
\bibitem{vanschoren2014} J.~Vanschoren, J.~N. van Rijn, B.~Bischl, and
L.~Torgo, ``OpenML: Networked science in machine learning,'' \emph{SIGKDD
Explorations}, vol.~15, no.~2, pp.~49--60, 2014.
\end{thebibliography}

\end{document}
